\chapter{Conclusões e recomendações gerais}


\section{Conclusões}

Com base no diagnóstico energético realizado, constatou-se que há um \textbf{potencial significativo de redução do consumo de energia elétrica} na unidade consumidora, especialmente nos sistemas de \textbf{\textcolor{red}{[iluminação / climatização / correção de fator de potência / outros]}}.

O modelo tarifário vigente, do tipo <<modalidadeTarifaria>>, associado à demanda contratada de \textbf{\textcolor{red}{[\_\_\_\_]}}~kW, mostrou-se \textbf{\textcolor{red}{[adequado / superdimensionado / subdimensionado]}} ao perfil de carga observado. Identificou-se, contudo, a possibilidade de \textbf{otimização contratual e técnica}, como a reavaliação da demanda, melhoria do fator de potência e racionalização do uso nos horários de maior custo.

As propostas de eficientização consolidadas no Capítulo~\ref{sec:acoes-propostas} apresentaram \textbf{viabilidade técnica e econômica}, com potencial de economia de \textbf{\textcolor{red}{[\_\_\_\_]}}~MWh/ano, redução de demanda de \textbf{\textcolor{red}{[\_\_\_\_]}}~kW e retorno financeiro estimado da ordem de \textbf{R\$~\textcolor{red}{[\_\_\_\_]}} por ano, reforçando a atratividade do projeto.

\section{Recomendações}

Considerando os resultados obtidos, recomenda-se a adoção das seguintes medidas de eficiência energética, conforme detalhamento apresentado no Capítulo~\ref{sec:acoes-propostas}:

\begin{itemize}
    \item \textbf{Substituir o sistema de iluminação:} uso de luminárias LED com sensores de presença e controle automatizado para ambientes de uso intermitente.

    \item \textbf{Modernizar a climatização:} substituição de equipamentos obsoletos por modelos com Selo Procel de Desempenho Energético e incorporação de sistemas de controle inteligente.

    \item \textbf{Aproveitar soluções passivas:} manutenção e uso de elementos arquitetônicos (claraboias, venezianas, exaustores) para reduzir o uso de iluminação e ventilação artificial.

    \item \textbf{Instalar sistema de monitoramento energético:} viabilizando controle em tempo real e diagnósticos contínuos para ações corretivas e preventivas.

    \item \textbf{Revisar ou religar banco de capacitores:} correção do fator de potência para evitar cobranças por energia reativa excedente.

    \item \textbf{Avaliar o modelo tarifário:} a partir do perfil de carga e das oportunidades de adequação contratual, visando a redução dos custos com demanda e energia.

    \item \textbf{Buscar fontes de financiamento:} por meio do \textit{Programa de Eficiência Energética da ANEEL (PEE)}, contratos de performance, ou recursos próprios, conforme a estrutura e porte da proposta.
\end{itemize}

A implementação dessas recomendações contribuirá diretamente para a \textbf{redução dos custos operacionais}, \textbf{melhoria do desempenho energético} e \textbf{sustentabilidade institucional}, promovendo maior controle e eficiência no uso dos recursos públicos.

